%==============================================================================
% UQFF Manuscript v2 — Section 8
% Topic:    Limitations, open questions, honest acknowledgments
% Status:   first draft, ready for Daniel's review.
% This is the section reviewers will read first to test honesty. Calibration matters.
% Anchors:  COVERAGE_GAPS.md, PREDICTION_LABELS.md, STATISTICAL_HYGIENE.md §7,
%           PAPER_1005 erratum, formal/UQFF/Millennium.lean (sorry policy)
%==============================================================================

\section{Limitations and open questions}
\label{sec:limitations}

A 1{,}994-whitepaper corpus reporting agreement across $\sim 800$
closures owes its readers a candid accounting of where the framework
falls short. This section enumerates the limitations a reviewer would
reasonably want explicitly disclosed.

\subsection{The closures are not formal proofs in the Clay sense}
\label{sec:lim_clay}

For the eight closures adjacent to Clay Mathematics Institute
Millennium Prize problems (Yang-Mills, Riemann, Navier-Stokes,
Hodge, Poincar\'e, P vs NP, Birch--Swinnerton-Dyer, and the
black-hole information paradox), the UQFF closures are
\emph{structural numerical identities matching the conjectured
value or anchor}. They are not formal proofs in the Clay sense.
The Clay statements require constructive existence proofs in the
Osterwalder--Schrader (Yang-Mills) or analytic-number-theory
(Riemann) or differential-geometry (Hodge) frameworks; UQFF
supplies numerical values, not constructive proofs of the
underlying mathematical statements.

This boundary is encoded explicitly in the framework's own Lean 4
scaffold (\texttt{formal/UQFF/Millennium.lean}, bundled with the
package), in a comment that reads ``the scaffold makes no assertion
that the UQFF claim implies the Clay statement.'' Removing the
\texttt{sorry} markers in that scaffold without independent peer
review at a Clay Qualifying Outlet would not constitute a Clay
proof, by the framework's own internal standard.

\subsection{No independent third-party reproduction}
\label{sec:lim_reproduction}

As of this writing, none of the framework's closures have been
independently reproduced by a third party. The package is
open-source and the reproduction infrastructure
(\texttt{pip install uqff}, fidelity gate, CI/CD, REST API,
Jupyter integration, C++ cross-language verification) is in place
to enable third-party reproduction; but the manuscript reports
single-author single-toolchain results, and the reader should
weigh the claims accordingly until independent groups confirm.

This is the most important single limitation. The Bonferroni and
$\Delta\mathrm{BIC}$ analyses of Sec.~\ref{sec:statistics} show
that the closure set is not a multiple-comparisons artifact, but
they cannot rule out a systematic bias common to a single
implementation. Independent re-implementation in a different
language (Julia, Rust, Fortran) by an independent group is the
single highest-value follow-up activity for the framework.

\subsection{Five tension closures with residuals above 1\,\%}
\label{sec:lim_tensions}

Five of the 263 schema-tagged closures carry the
\textsc{prod\_tension\_or\_outlier} classification:

\begin{table}[H]
\centering
\caption{The five \textsc{prod\_tension\_or\_outlier} closures with
residuals above 1\,\%. None are silently failing; each is documented
in \texttt{forward\_predictions.md} as either a deliberate
disagreement with one of two competing measurements or a flagged
refinement target.}
\label{tab:tensions}
\begin{tabular}{l r p{6.5cm}}
\toprule
\textbf{Closure} & \textbf{Residual} & \textbf{Disposition} \\
\midrule
$\sin^2\theta_W$ (electroweak) & 3.4\,\% & Renormalization-group running unmodeled at current closure depth; PAPER\_1209HH \\
$\alpha_s(M_Z)$ (strong coupling) & 13.7\,\% & Scheme dependence ($\overline{\mathrm{MS}}$ vs.\ on-shell); refinement target \\
Pons--Fleischmann Pd-D excess power & $\sim 100\times$ overshoot & Reactor-geometry under-constrained; closure machinery applied without parameter retuning \\
Rossi E-Cat COP family & $\sim 3$--$5\times$ undershoot & Same disposition as Pons--Fleischmann (see Sec.~\ref{sec:holmlid}) \\
Star-Magic reactor COP & $\sim 30\times$ undershoot & Either closure underweights a Star-Magic-specific factor, or the reported COP requires independent reproduction (Sec.~\ref{sec:holmlid}) \\
\bottomrule
\end{tabular}
\end{table}

We do not paper over these. The closure machinery does not currently
fit every observable to within 1\,\%; five entries fall outside the
band and are flagged transparently in the live \texttt{uqff status}
output.

\subsection{The Yang-Mills 1.736 GeV registry-bug history}
\label{sec:lim_ym_history}

Between 2026-04 and 2026-06-25, the calculator's Yang-Mills
dispatcher returned a hardcoded \texttt{1.736} GeV value that did
not match the formula in the cited source paper (PAPER\_1005) and
did not match the lattice QCD anchor at $\sim 1.7$\,GeV. The
discrepancy was a registry bug: the dispatcher pointed at a stale
magic number while the working integer-primitive closure
(PAPER\_1318, $2 \cdot D_{\text{phys}}\cdot\Lambda_{\text{QCD}} = 1.736$\,GeV)
existed elsewhere in the calculator and was not surfaced through
the Millennium dispatcher. The bug was identified and corrected on
2026-06-25 (see Sec.~\ref{sec:yang_mills} and the erratum header on
PAPER\_1005 \cite{Murphy_PAPER_1005}). All 610 propagated citations
of ``1.736 GeV (PAPER\_1005)'' in the GW-bucket whitepapers were
updated in-place to cite PAPER\_1318 with the corrected value.

This is the most embarrassing single item in the framework's history.
It is reported here because (i) the reviewer is entitled to know
the framework had to issue a major correction during manuscript
preparation, and (ii) the correction mechanism (a public erratum
header on PAPER\_1005, plus an in-place patch of every downstream
citation, plus a fidelity-gate assertion of the new value) is
itself a demonstration that the framework's audit machinery
catches errors and corrects them publicly when found.

\subsection{Two neutrino mass-squared splittings are anchored, not predicted}
\label{sec:lim_neutrinos}

Table~\ref{tab:sm_fermions} reports the two neutrino mass-squared
splittings ($\Delta m^2_{21} = 7.5\times 10^{-5}$\,eV$^2$,
$\Delta m^2_{31} = 2.5\times 10^{-3}$\,eV$^2$) at residual
$0.000\,\%$ against the global-fit central values
\cite{Esteban2020}. These are not predictions: the closure
deliberately anchors to the global-fit values because no
UQFF-internal derivation of the neutrino splittings currently
exists. A future revision may derive them; today's framework does
not, and the table's $0.000\,\%$ entry should be read accordingly.

\subsection{1{,}396 whitepapers without a direct dispatch-key backing}
\label{sec:lim_orphan_whitepapers}

Of the 1{,}994 whitepapers bundled with the package
(Sec.~\ref{sec:whitepaper_corpus}), approximately 1{,}396 are not
directly referenced by any closure's \texttt{primary\_source}
attribute. These are not orphan or junk papers --- they cover
historical derivation chains, alternative formulations, audit
artifacts, and exploratory material --- but the parameter-economy
claim of Sec.~\ref{sec:bic} is anchored on the 598 whitepapers
that \emph{do} back specific closures; the additional 1{,}396 are
context, not load-bearing.

A full whitepaper-to-closure mapping audit is provided in
\texttt{COVERAGE\_GAPS.md} (bundled). Future work would either
(i) tighten the mapping to bring the inverse-coverage rate higher,
or (ii) document a triage of which whitepapers are
foundational versus exploratory, so the reader can navigate the
corpus by relevance rather than by paper number.

\subsection{Three primitives may yet prove derivative}
\label{sec:lim_primitives}

As discussed in Sec.~\ref{sec:open_provenance}, three of the
quantities currently treated as locked
($\mathrm{SSq} = 0.57$, $S_{26}$, $\omega_{\text{SCm}} = 1.25\,\mathrm{THz}$)
carry a C-grade provenance in
\texttt{PROVENANCE\_AUDIT.md} and may be reducible to lower-grade
primitives via further analysis. If any of the three prove
derivative, the true free-parameter count drops below 9, the BIC
argument strengthens, and the closure structure simplifies. The
conservative posture taken in this manuscript is to assume nine
primitives until the reductions are formally established.

\subsection{The Star-Magic reactor experimental claim is unconfirmed}
\label{sec:lim_star_magic}

The Star-Magic reactor architecture cited in Sec.~\ref{sec:holmlid}
(Table~\ref{tab:lenr_reactors}, last row) has been reported by the
framework's author at COP 555 at 27\,W input, ambient temperature,
$\mathrm{pH}-37$. Independent reproduction of this performance has
not been published. The framework's own closure machinery does not
reproduce the reported COP (UQFF predicts $\sim 18$, observed 555,
a $\sim 30 \times$ undershoot). The honest accounting is that
either (a) the closure machinery is underweighting a geometric or
thermodynamic factor specific to the Star-Magic geometry, or
(b) the reported COP itself requires independent confirmation, or
both. We do not adjudicate between (a) and (b) in this manuscript;
both are flagged as open work.

\subsection{Manuscript versioning note}
\label{sec:lim_manuscript}

This manuscript (\texttt{manuscript\_v2/}) is the
parameter-economy physics paper targeting Foundations of Physics
peer review. An earlier independent draft
(\texttt{Manuscript 1\_12Feb2026/}, also bundled with the package)
is a software-implementation paper covering REST API performance,
production-deployment patterns, and computational architecture;
that draft targets a different venue (likely SoftwareX or Computer
Physics Communications) and is structurally distinct from the
current document. Both manuscripts exist in the repository; the
present submission is the v2 physics paper.

\subsection{Closure-runner target-dict calibration carry-overs}
\label{sec:lim_runner_targets}

The \texttt{run\_millennium\_proofs.py} runner reports two of seven
Millennium closures (\texttt{navier\_stokes} and \texttt{poincare})
as having large residuals against the runner's internal target dict
because the units or scoring conventions used in the target dict
differ from the UQFF closure's natural output. These are
calibration mismatches in the runner script, not closure failures;
fixing them is queued for v5.30.0 but does not affect any value
reported in this manuscript.

\subsection{Summary of limitations}
\label{sec:lim_summary}

Three limitations dominate the framework's current epistemic state:

\begin{enumerate}
  \item \textbf{No independent reproduction yet.} Single-author,
        single-toolchain results. Independent re-implementation is
        the highest-value follow-up.
  \item \textbf{Closures are numerical identities, not Clay-style
        proofs.} The Lean 4 scaffold makes this explicit; the
        manuscript repeats it in every Clay-adjacent section
        (Secs.~\ref{sec:yang_mills} and \ref{sec:lim_clay}).
  \item \textbf{Five tensions remain disclosed but unresolved.}
        Section~\ref{sec:lim_tensions} enumerates them; none are
        suppressed.
\end{enumerate}

The framework explicitly does not claim:

\begin{itemize}
  \item to be a proven theory of everything,
  \item to replace the Standard Model or $\Lambda$-CDM (the
        ``NOT REPLACEMENT'' clause runs through every UQFF
        whitepaper),
  \item to have proven any Clay Millennium prize problem,
  \item to have eliminated all parameter freedom (nine primitives
        remain),
  \item to have been validated by independent experimental
        reproduction.
\end{itemize}

What the framework does claim --- $\Delta\mathrm{BIC} = 94.1$ on
parameter-economy grounds, 226 Bonferroni-passing closures across
253 observations, 42 falsifiable forward predictions --- is an
invitation for independent scrutiny, not a final verdict.

